\chapter{Introduction}\label{ch:intro}

A vehicle drives through a city it has never visited and meets a truck carrying an oddly shaped load that overhangs its bed. A robot that has learned to fold shirts from hours of human demonstrations is handed a garment crumpled in a way that no demonstrator ever produced. Neither situation is exotic; each happens every day somewhere. Yet each is, for the agent concerned, a condition that its training data did not contain. This thesis is about such conditions, and about how physical agents can be built to cope with them.

% ==============================================================================
\section{Physical Agents and the Curse of Reality}\label{sec:intro:curse}

A \emph{physical agent} perceives the world through sensors, reasons about what it perceives, and acts on the world through actuators, in a closed loop. Autonomous vehicles and robot manipulators are the two kinds of physical agents studied in this thesis. Over the past decade, both have increasingly been built by learning from data: perception models are trained on annotated sensor recordings, planners and control policies are trained to imitate human demonstrations, and, most recently, foundation models pre-trained on internet-scale data have been adapted to driving and manipulation.

Learning from data brings a fundamental difficulty with it. A learned model is reliable only on data that resemble its training data, whereas a physical agent is deployed in an open world. Rare events are individually improbable, but so numerous that, taken together, they occur all the time; for autonomous vehicles this has been called the \emph{curse of rarity}~\cite{liu2024curse}. In addition, a physical agent acts in closed loop: its own actions determine the situations it encounters next, so a policy that has only seen an expert's behaviour drifts into states the expert never visited, and its errors compound~\cite{ross2011dagger}. This thesis refers to the combination of these difficulties as the \emph{curse of reality}: whatever data an agent is trained on, the real world will keep presenting conditions the agent has never seen.

The curse of reality is more severe for physical agents than for most other applications of machine learning, for four reasons. Their mistakes have physical consequences, so failures on rare inputs cannot be dismissed as statistically negligible. Their data are expensive, because every recording requires a vehicle or a robot and every annotation of 3D structure or of a manipulation trajectory requires skilled effort. Their loop is closed, so small errors do not stay small. And their deployment platforms impose hard limits on latency, memory and computation, so that robustness cannot be bought simply by running larger models.

% ==============================================================================
\section{From Modular Pipelines to End-to-End Physical Agents}\label{sec:intro:context}

The research in this thesis was carried out between 2021 and 2026, a period in which the design of physical agents changed substantially. Autonomous driving systems were traditionally built as modular pipelines, in which separately trained perception, prediction and planning modules communicate through hand-designed interfaces such as lists of 3D bounding boxes. Two developments challenged this design. The first was the rise of unified scene representations learned directly from multiple cameras, such as the bird's-eye-view (BEV) representation of BEVFormer~\cite{li2022bevformer}, which made it possible to train one network for many perception tasks and to share features across them~\cite{li2023bevdevils}. The second was end-to-end planning: UniAD showed that when all modules are trained jointly towards the final planning objective, planning error and collisions drop substantially~\cite{hu2023uniad}.

A third development followed: foundation models. Vision-language models (VLMs) trained on web-scale data carry broad knowledge about the world and can express their reasoning in language, which raised the prospect of drivers that can explain their decisions and generalise to unfamiliar situations~\cite{sima2024drivelm}. In robotics, vision-language-action models and large, openly released robot datasets promised generalist manipulation policies~\cite{bu2025agibot_iros}. Each of these developments helps against the curse of reality in some respects and introduces new risks in others: unified representations must still describe objects they were not designed for, end-to-end planners are hard to guard with rules, language outputs can look competent without being grounded~\cite{xie2025vlms}, and large policies trained on large data still fail when training and deployment conditions diverge~\cite{sima2026kai0}. This thesis studies these risks and proposes ways to address them.

% ==============================================================================
\section{Research Questions}\label{sec:intro:rq}

The central question of this thesis is: \emph{how can physical agents be made robust to the conditions that the real world presents but their training data do not contain?} Following the perceive--reason--act loop, it is divided into four research questions.

\begin{description}[leftmargin=1.2cm, labelwidth=1cm, itemsep=4pt]
  \item[RQ1] \textbf{Perceive.} How should a physical agent represent its 3D surroundings so that the representation remains valid for the open set of objects and structures met in the real world, while remaining learnable from cameras and useful for planning?
  \item[RQ2] \textbf{Reason.} How can the general knowledge of vision-language models be grounded in driving in a structured and verifiable way that mirrors how humans reason from perception through prediction to planning, and does such reasoning help a driving agent generalise to unseen situations?
  \item[RQ3] \textbf{Act.} How can a learned policy remain reliable when the deployment distribution departs from its training distribution: (a) for end-to-end driving, without hand-designed rules, cost functions or additional labels; and (b) for real-world robot manipulation, under inconsistent demonstrations and a limited budget of data and computation?
  \item[RQ4] \textbf{Architecture.} What system-level design allows perception, reasoning and action to correct one another, so that a physical agent can keep improving after deployment?
\end{description}

RQ1 to RQ3 are each answered by one or two core chapters. RQ4 is a synthesis question: it is answered in Chapter~\ref{ch:discussion} by drawing on the evidence of all core chapters and of the author's co-authored work.

% ==============================================================================
\section{Approach}\label{sec:intro:approach}

The research follows a benchmark-driven methodology, described in detail in Chapter~\ref{ch:methodology}. Each project begins with a failure observed, or anticipated, in the real world. The failure is formalised as a task with a representation and a metric, turned into a public benchmark, and addressed with a simple and scalable baseline method; the benchmark is then opened to the community, often through a public challenge, and the method is finally tested in closed loop or in the real world, where it is made to correct itself under the shift it meets. The residual failures found at that point start the next iteration. Fig.~\ref{fig:intro:overview} shows how the core chapters of the thesis map onto the perceive--reason--act loop and the research questions.

\begin{figure}[t]
  \centering
  \tikzfig{ch-introduction/figures/thesis_overview}
  \caption[Overview of the thesis]{Overview of the thesis. The core chapters follow the perceive--reason--act loop of a physical agent facing the curse of reality: Chapter~\ref{ch:occnet} addresses perception (RQ1), Chapter~\ref{ch:drivelm} reasoning (RQ2), and Chapters~\ref{ch:centaur} and~\ref{ch:kai0} acting in driving and manipulation (RQ3). Chapter~\ref{ch:discussion} synthesises them into an answer to the architectural question RQ4. Chapter~\ref{ch:methodology} describes the research methodology that underlies all of them.}
  \label{fig:intro:overview}
\end{figure}

% ==============================================================================
\section{Contributions}\label{sec:intro:contributions}

The main contributions of this thesis are the following.

\begin{enumerate}[label=\textbf{C\arabic*}, leftmargin=1.1cm, itemsep=4pt]
  \item \textbf{Occupancy as a universal scene representation (Chapter~\ref{ch:occnet}).} We propose to describe driving scenes by 3D semantic occupancy instead of 3D boxes, together with OccNet, a camera-based network that reconstructs occupancy through a cascade of voxel decoders, and OpenOcc, a dense occupancy benchmark with 34,149 frames and more than 1.4 billion annotated voxels. Occupancy proves to be a general descriptor: it improves semantic scene completion (26.98 mIoU against 23.67 for the previous camera-based state of the art), describes irregular objects better than boxes, serves as a pre-training target for detection, and lowers the collision rate of a planner by 15--22\% compared with planning on 3D boxes, and by up to 58\% compared with planning on BEV segmentation. We further scale occupancy supervision with OpenScene, a compact redistribution of more than 120 hours of nuPlan driving data with occupancy labels, which has served as the data of the end-to-end driving and world-model tracks of the CVPR 2024 and 2025 challenges.
  \item \textbf{Structured reasoning with vision-language models (Chapter~\ref{ch:drivelm}).} We introduce Graph Visual Question Answering, in which perception, prediction, planning, behaviour and motion questions are linked by logical dependencies, together with the DriveLM-nuScenes and DriveLM-CARLA datasets (4,871 frames with 91.4 question--answer pairs per frame, and about 1.6 million question--answer pairs, respectively) and the DriveLM-Agent baseline. We show that graph-structured reasoning improves zero-shot generalisation to a new sensor configuration and to unseen object categories. The benchmark underpinned the ``Driving with Language'' track of the CVPR 2024 Autonomous Grand Challenge, with 152 participating teams.
  \item \textbf{Test-time self-correction for end-to-end driving (Chapter~\ref{ch:centaur}).} We propose Centaur, which adapts a deployed planner by test-time training on Cluster Entropy, a new uncertainty measure that requires neither labels nor hand-designed rules. Centaur raises the driving score on NAVSIM from 90.3 to 92.6 and ranked first among three-seed entries on the NAVSIM leaderboard at the time of submission. We also introduce \texttt{navsafe}, a split of safety-critical scenarios that exposes the remaining gap to human driving.
  \item \textbf{Robust manipulation under distributional inconsistency (Chapter~\ref{ch:kai0}).} We analyse the failures of real-world manipulation policies as inconsistencies between the distributions of demonstrations, of what the policy can represent, and of what is executed on hardware, and propose \chiz, which addresses them with Model Arithmetic, Stage Advantage and Train--Deploy Alignment. With about 20 hours of data per task and eight A100 GPUs, \chiz surpasses the success rate of the state-of-the-art $\pi_{0.5}$ policy by nearly 250\% on garment manipulation and operates for 24 hours without interruption.
  \item \textbf{A benchmark-driven methodology and a synthesis (Chapters~\ref{ch:methodology} and~\ref{ch:discussion}).} Drawing on all eighteen publications of the author, we formulate a methodology for research on physical agents, consisting of a taxonomy of the shifts that make up the curse of reality, a benchmark-driven research cycle, five design principles and a ladder of evaluation fidelity, and we synthesise the findings of the core chapters into lessons for self-correcting architectures.
\end{enumerate}

% ==============================================================================
\section{Thesis Outline}\label{sec:intro:outline}

The remainder of the thesis is organised as follows. Chapter~\ref{ch:background} reviews the literature on distribution shift for physical agents, 3D scene representations, end-to-end driving, language models in driving, adaptation methods, robot manipulation learning, and benchmarks, and identifies the gaps that motivate the research questions. Chapter~\ref{ch:methodology} presents the research methodology. Chapters~\ref{ch:occnet} to~\ref{ch:kai0} present the core contributions along the perceive--reason--act loop: occupancy perception, graph-structured reasoning, test-time training for driving, and robust manipulation. Chapter~\ref{ch:discussion} interprets the findings across chapters, answers the research questions, and discusses limitations. Chapter~\ref{ch:conclusion} concludes the thesis and outlines directions for future work. The appendices contain supplementary material for each core chapter.

% ==============================================================================
\section{Publications}\label{sec:intro:publications}

The core chapters are based on the following publications, which the author led: Chapter~\ref{ch:occnet} on~\cite{sima2023occnet,sima2023openscene}, Chapter~\ref{ch:drivelm} on~\cite{sima2024drivelm}, Chapter~\ref{ch:centaur} on~\cite{sima2024centaur}, and Chapter~\ref{ch:kai0} on~\cite{sima2026kai0}. The co-authored works that inform Chapters~\ref{ch:methodology} and~\ref{ch:discussion} are~\cite{li2022bevformer,chen2022persformer,li2023bevdevils,hu2023uniad,gao2023sdf,wang2023openlanev2,huang2023vcd,ding2024hint,xie2025vlms,hamdan2025eta,bu2025agibot_iros,chen2025robot}, and the author's earlier work on efficient simulation and learning is~\cite{sima2021lshsmile}. The complete list, with the author's role in each work, is given in the List of Publications at the beginning of this thesis.
